\clearpage
\section{The full copying template and contextual improvement}
\label{sec:copying-full-copying-template}
The sixth of the seven further template families is now complete. The
literal copying pre-template has output closure, input safety at every
ordered input, universal forward and converse robustness, contextual
improvement and termination equivalence. All observing contexts in the
original unrestricted class are covered, including recursive occurrences
under boxes. Parametricity is the only one of those seven families whose
universal template proof remains open. Proposition~12.2 also still requires
the correspondence, semantic weakening and final Figure~22 assembly.

\subsection{The literal equation and its costs}
\termdef{full-copying-template}
Let $T$ be any genuine term contraction tree with $n+1$ ordered inputs,
and let $H:\mathsf{term}\Rightarrow\mathsf{term}^k\mathsf{atom}^h$ be an
original copyable, focus-free payload. Write $T^{-}$ for the original
residual tree: feed input zero of $T$ with an actual weakening, leaving its
other $n$ inputs in order. The token-erasure of Figure~17 compares
\[
 T;H\quad\sim\quad
 \operatorname{pair}_{k,h}(H,T^{-};H);
       (D^{\mathsf{term}}_{k,2}\otimes D^{\mathsf{atom}}_{h,2}).
\]
Here the pair constructor keeps the active input first, then the residual
inputs, and places the two term-resource blocks before the two atom-resource
blocks. Both distributors are the original canonical graphs. In particular,
their weakening edges are retained. The pair relation permits independent
recursive carrier names on both sides, fixes every external position, and
contains no semantic or robustness premise.

\termdef{copying-contextual-improvement}
The conclusion is quantitative improvement, not equality of exact costs.
For every successful left run of length $p$, there is a successful right
run of length $q\leq p$. Conversely every successful right run has a
successful left run whose length is at least the right run's length.
Forgetting these inequalities gives termination equivalence. This says
nothing that assumes a callback terminates: each implication is conditional
on the successful run it transports, while stuckness remains distinct from
successful termination.

The distinguished input explains the inequality. The left graph requests
copying and performs its original rewrite in exactly two transitions; it
then is the right graph with the same search token. The right side takes
zero transitions. At every other input, both sides take an actual
request/rewrite pair, with costs two and two. The resulting three payloads,
remaining tree and nested canonical resource rows have a proved sharing
interchange. These establish the original input-safety clauses. Applying
the exact cost preorder here would incorrectly require $2=0$; the published
contextual theorem instead uses the proved improvement and converse-overhead
cost triples.


\lean{Copying/CopyingTemplateInputSafety}{63}{65}{CopyingTemplate.template_inputSafe_improvement}
\par\penalty0
\lean{Copying/CopyingTemplateEquivalence}{16}{22}{CopyingTemplate.filling_improvement}
\par\penalty0
\lean{Copying/CopyingTemplateEquivalence}{27}{33}{CopyingTemplate.filling_reverse_overhead}
\par\penalty0
\lean{Copying/CopyingTemplateEquivalence}{37}{42}{CopyingTemplate.filling_equiv}
\par\penalty0

\subsection{Classifying the original remote copying matches}
The eight operation-headed local-rule families have inspected primitive
footprints, original stable arguments and, for beta, the actual lambda and
its complete box body. A shallow left copying component starts with a term
contraction, so it cannot supply one of these incompatible primitive heads.
The forward ancestry barrier retains every inspected edge. In the converse,
original rooted readiness and wire typing exclude an incompatible shallow
passive contribution to an operation's evaluated argument. Minimal atom
spines retain exactly the store/equality entries inspected by the local
rule; unused siblings stay attached in the common complement. Thus actual
matching rewrites and original residuals exist for all eight operation cases.

The copying-headed case requires an additional diagram.
\termdef{copying-first-shallow-entry}
Follow the actual finite contraction resolution from the original token to
its complete copyable payload. Either the selected route and payload are
retained, or there is a first shallow copying hole, a specific ordered input
$j$ of that hole, a retained prefix, and a suffix ending in that hole's actual
payload. Functional resolution identifies the complete payload, including
its recursive box label; it does not identify only its passive head.

In the retained case a minimal selected term spine replaces the arbitrary
whole tree in the supplied match. Its active input and payload attachment
are unchanged. All inspected edges avoid shallow holes, so the proved
retained-footprint construction transfers the actual source occurrence.
Successor coherence compares the rematched rewrite with the independently
supplied original rewrite. The unused branches are retained in the host;
no condition is imposed on their labels or recursive fillings.

\termdef{copying-input-graft}
In the shallow case let $P$ be the retained prefix tree. Grafting $P$ into
input $j$ of $T$ uses a literal two-component frame with one new junction.
The upstream inputs come first, followed by every unselected downstream
input in its original relative order. The construction allows a zero-input
upstream residual tree. It adds no identity edge and takes no quotient of
input order. Its raised primitive-tree presentation has a proved recursive
isomorphism to the actual two-component hypernet.

The complete original prefix occurrence and the original hole corolla are
edge-disjoint: sharing labels cannot equal the hole label. Complete incidence
then supplies their internal freshness. Joining them gives an actual
occurrence of this graft in the original common context. Label-copy
certificates reconstruct its two actual filled sides. A finite common cut
lifts any proved literal diagram through this occurrence, preserving the
original token, all other hole names and every untouched recursive filling.


\lean{Copying/ContractionTreeGraft}{100}{116}{ContractionTree.graftInput}
\par\penalty0
\lean{Copying/ContractionTreeGraftRewrite}{52}{62}{ContractionTree.Graft.copy_source_iso}
\par\penalty0
\lean{Copying/CopyingRemoteCut}{35}{47}{CopyingTemplate.RemoteCut.diagram}
\par\penalty0

\subsection{Positive and distinguished shallow inputs}
\termdef{copying-positive-input-diagram}
If $j>0$, input zero of $T$ is still an external port of the full graft,
separate from the upstream active port. The original rewrite confined to
$T;H$ and a rewrite using the full graft tree start in exactly the same whole
graph at that same external port. Proved successor coherence therefore
identifies their actual reducts. After reindexing the two selected original
inputs, the already proved two-input copying interchange constructs the
required residual. Undoing the actual input permutation restores every
ordered external position. No new copying commutation law is assumed.

\termdef{copying-distinguished-input-diagram}
If $j=0$, the downstream selected port is internal to the graft. The
positive-input argument does not apply. First use the independently proved
sharing algebra to relate the actual residual of the full graft to the
graft of $P^{-}$ into $T$. A fixed-boundary partial-serial reassociation
isomorphism then exposes one genuine original copying-template instance
$T;H$ inside that retained branch. Applying that original relation leaves
three actual branches
\[
 H,\qquad P^{-};H,\qquad T^{-};H.
\]
Each original resource now has two complete canonical two-use rows.

\termdef{copying-three-branch-merging}
Merging three arbitrary focus-free branch graphs through these rows is
associative up to a genuine sharing-algebra filling. The proof gives a
common finite frame with one hole for each entire branch and one three-input
sharing hole for each original term or atom resource. An explicit input
bijection fixes the order of all three branch input blocks, including empty
blocks. Each association has complete original branch and row occurrences
in that frame, exact attachment equations, and pairwise edge disjointness.
An exhaustive partition covers every actual edge of each association.
The complete occurrence and fixed external boundaries therefore yield a
recursive graph isomorphism. The two six-edge rows are genuine same-input
sharing trees, so the sharing template supplies the relation between them.
Arbitrary branch graphs are reflexive contributions: no branch is executed,
simplified or identified with another branch.

On the right side of the distinguished diagram, factor the actual graft
through the active branch of the original copying reduct. This yields the
literal merge of $P;H$ and $T^{-};H$. One original copying rewrite in its
first branch gives the other association of the same three branches.
The proved merging law supplies the final sharing correction. All these
are whole-graph constructions; a local resource-row picture alone would
not establish the required original residual specimen.


\lean{Copying/CopyingMergeAssociativity}{11}{20}{CopyingMerge.assoc}
\par\penalty0
\lean{Copying/CopyingRemotePositiveRaw}{15}{33}{CopyingRemote.Positive.raw_diagram}
\par\penalty0
\lean{Copying/CopyingRemoteZero}{88}{106}{CopyingRemote.Zero.raw_diagram}
\par\penalty0
\lean{Copying/CopyingRemotePair}{33}{40}{CopyingTemplate.Pair.graft_diagram}
\par\penalty0

\subsection{Two-sided corrections and universal robustness}
\termdef{copying-two-sided-residual}
The distinguished case records a leading sharing witness, one original
copying witness and a trailing sharing witness. This is not transitivity
of the original copying filling relation. At each link the actual graph,
search token and ordered interface are recorded. Expansion acts reflexively
on unrelated components in the sharing links and retains their original
copying certificates in the middle link.

The independently proved exact sharing semantics supplies the two
refinements around an original quasi-specimen. Its required intermediate
roots are derived from the actual rooted endpoints by sharing's nonrewrite
root-reflection theorem. Hole freedom is transported by the actual label
laws. This discharges the original refinement-up-to clause without an
assumed residual-admissibility certificate.

Both literal cases preserve the upstream input zero. Graft congruence
transports their actual rewrites through independent recursive graph names.
Their union therefore proves the remote diagram for every original copying
pair and every ordered hole input. The common-cut construction turns it into
actual matching whole-graph rewrites and a genuine original residual.

The converse classification also includes the right graph's zero-hop
resolution at input zero and its positive-input residual-tree resolutions.
Together with retained rematching and the eight operation-headed cases,
these exhaust all nine original local-rule constructors in both directions.
Reflection of an arbitrary representation-class successor uses the existing
proved machine laws. Consequently the universal forward and converse
robustness theorems have no matching, classification, determinism,
refocusing or residual-existence hypotheses left to discharge.

Concrete sufficiency, the proved input-safety inequality and original initial
lifting give contextual improvement. Converse sufficiency gives the reverse
overhead comparison. Costs are forgotten only when deriving termination
equivalence. The result covers arbitrary original addressed sites under
boxes as well as ordinary closing contexts. Pointed variants retain their
explicit actual-root premises. The original terms, Figure~16 translation
and Figure~17 operational rule are unchanged.


\lean{Copying/FullCopyingForwardRobustness}{101}{126}{CopyingTemplate.template_robust}
\par\penalty0
\lean{Copying/CopyingTemplateConverseRobustness}{120}{149}{CopyingTemplate.template_robust_converse}
\par\penalty0
\lean{Copying/CopyingTemplateEquivalence}{61}{66}{CopyingTemplate.addressed_improvement}
\par\penalty0
\lean{Copying/CopyingTemplateEquivalence}{70}{75}{CopyingTemplate.addressed_equiv}
\par\penalty0
\lean{Copying/CopyingTemplateEquivalence}{143}{150}{CopyingTemplate.pair_contextual_equiv}
\par\penalty0

\subsection{Checked parametricity progress, with universal work still open}
\termdef{parametricity-private-read-correction}
The actual private-reference beta/read/copy path has a checked nine-transition
prefix with an arbitrary original stable argument. Its exact computed
successor has a genuine copying-template filling followed by two independent
garbage-collection fillings. Complete occurrences retain the store and all
remaining reference inputs. A full edge-cover proof reconstructs a returned
numeral one together with the surviving tree, original store and stored
numeral one. Lifting the correction through the original computed complement
keeps every unused argument output. These are structural corrections with
proved quantitative semantics; they are not extra machine transitions.

\termdef{parametricity-large-beta-opening}
For the larger abstraction, the two original beta rules each take one
transition. Their actual reducts are recursively isomorphic to applying the
corresponding smaller abstraction to the original callback graph connected
to the evaluated argument. The callback retains its literal single-use
contraction tree and unit edge. A four-component complete occurrence
accounts for the application head, smaller abstraction, callback and argument;
edge coverage and fixed boundaries prove the isomorphism. This does not
execute the callback and has no callback-termination hypothesis.

The remaining family isolation, ordered reattachment, shallow lambda
copying and universal parametricity robustness are separate obligations.
Neither a counted prefix nor these local opening diagrams prove the
parametricity template. Appendix~B.1 keeps those tasks visible.

\lean{Copying/ParametricityPrivateCorrection}{115}{125}{Parametricity.PrivateApplication.nine_transition_correction}
\par\penalty0


\lean{Copying/ParametricityLargeOpening}{49}{51}{Parametricity.LargeOpening.read_beta}
\par\penalty0


\lean{Copying/ParametricityLargeOpening}{55}{57}{Parametricity.LargeOpening.one_beta}
\par\penalty0

